\section{Hermitian unitals}
Let $K=\mathbb F_{q^2}$ and let $V$ be a nondegenerate three-dimensional
Hermitian space. Its isotropic one-spaces, with the intersections of
projective secants as blocks, form the classical
$2$-$(q^3+1,q+1,1)$ unital $U_H(q)$. Each point belongs to $q^2$ blocks.
These standard design parameters are the same ones used in [L].

\textbf{Theorem H.Unital.}
For every prime power $q\geq2$, $\Gamma(U_H(q))$ is firm and residually
connected and satisfies $(\mathrm{IP})_2$. The group $PGU(3,q)$ acts
regularly on its chambers.

\emph{Proof.}
The parameter $r=q^2$ satisfies the hypotheses of [G, Theorems S.1 and
S.6], proving firmness and the rank-two intersection property. Directly,
the three panel sizes for point--block, point--flag and block--flag pairs
are $q+1$, $q^2-1$ and $q$.
At a flag $(p,M)$ the residue is $K_{q,q^2-1}$. At a point $Q$, its flag
objects are $(p,Qp)$ for $p\ne Q$ and its block objects are the blocks
avoiding $Q$; incidence is $p\in L$. If $p,r$ lie on different blocks
through $Q$, their joining block avoids $Q$ and connects them. If they
lie on the same such block, choose $s$ outside it and use $ps,sr$.
Every remaining block has point neighbors, so the point residue is
connected. The block residue is connected by [G, Theorem H.3], which
identifies it with the full isotropic-direction graph on a nonzero norm
shell. All field, transport and generation arguments for this step are
proved in that common paper.

The whole incidence graph is connected. Given distinct original points
$Q,R$, choose $p$ outside their block. At most two of the $q^2$ blocks
through $p$ contain $Q$ or $R$, leaving one incident in $\Gamma$ with
both. Every block has an outside point and every flag $(p,M)$ has a
neighbor in $M\setminus\{p\}$. These are all residues of rank at least
two, proving residual connectedness.

For the action, a chamber is $(Q,L,(p,M))$ with $Q\ne p$ on $M$ and
$p\in L\ne M$. The unitary group is transitive on ordered distinct
isotropic points: choose representatives $e,f$, rescale their nonzero
pairing to one, complete by a vector of norm one orthogonal to both,
and map one such basis to another. Fix $p=\langle e_1\rangle$ and
$Q=\langle e_3\rangle$ in coordinates with
\[
h(x,x)=x_1x_3^q+x_1^qx_3+x_2^{q+1}.
\]
All projective lines through $p$ other than $M=pQ$ have the form
$L_t=\langle e_1,e_2+t e_3\rangle$, $t\in K$.
The line $L_0$ is tangent. For $t\ne0$ the equation
$x t^q+x^qt+1=0$ has $q$ solutions, so $L_t$ is a secant.
The projective stabilizer of $(p,Q)$ is exactly
\[
D_a=\operatorname{diag}(a,1,a^{-q}),\qquad a\in K^*.
\]
Indeed any such isometry preserves the two isotropic lines and their
one-dimensional orthogonal complement. Its scalar on the latter has
norm one; dividing by that scalar gives the displayed matrix. It sends
$L_t$ to $L_{t a^{-q}}$, regularly on the $q^2-1$ possibilities. Thus
$PGU(3,q)$ is chamber-transitive with trivial chamber stabilizer. The
chamber count is
$q^3(q^3+1)(q^2-1)$, also obtained from [G, Theorem S.3]. \hfill$\square$

For $q>2$ we now recover the diameter and gonality data in Figure 2 of
[L]. For a bipartite incidence graph, $(d_0,g,d_1)$ denotes the two
vertex-type eccentricities and half its girth, using the source's
convention.
